\section{Solution Design}
\label{sec:design}

\subsection{Panoramica architetturale}
Il sistema è composto da $K$ nodi \emph{worker} paritari e da un
\emph{discovery server} leggero (Fig.~\ref{fig:architecture}). Ogni worker
esegue due thread indipendenti: un server gRPC che riceve i modelli inviati
dai vicini (Fase C di un altro worker) e li accumula in un buffer condiviso,
e un ciclo di training che, a ogni round, aggrega quanto ricevuto, esegue
addestramento locale e invia i propri pesi a un sottoinsieme di vicini.
Qualunque worker può avviare il proprio processo di training in autonomia,
leggendo iperparametri e riferimenti dei peer coinvolti da una configurazione
condivisa (\texttt{config.yaml}) e dal discovery server; non esiste un nodo
che comandi l'avvio degli altri. Questo soddisfa la decentralizzazione
dell'\emph{avvio} richiesta dalla traccia — nessun coordinatore privilegiato
è necessario. Il sistema resta però un job sincronizzato all'avvio, non un
servizio always-on: prima di iniziare il training, ogni worker attende
(Sezione~\ref{sec:details}A) che l'intera coorte di $K$ client configurati
sia online, con un timeout oltre il quale procede comunque con i peer
disponibili; un client non può quindi iniziare una nuova sessione in un
istante arbitrario senza un nuovo deployment, un limite condiviso con la
generalità dei job di training distribuito (analogamente a un job
SLURM/MPI) e discusso in Sezione~\ref{sec:conclusions} come possibile
estensione.

Il discovery server espone quattro sole operazioni REST
(\texttt{/register}, \texttt{/deregister}, \texttt{/peers}, \texttt{/health}),
mantiene in memoria una mappa \texttt{worker\_id~$\to$~indirizzo} e non
riceve, memorizza o media mai un peso del modello. Il suo unico compito
aggiuntivo è un \emph{watchdog} che forza la terminazione del cluster se un
worker non si deregistra correttamente entro un timeout — una funzione di
coordinamento della terminazione, non di training. Questo confina la
centralizzazione ai soli compiti di avvio/terminazione e identificazione dei
client, come richiesto dalla traccia. Il raggio d'azione di questo singolo
componente centralizzato è inoltre limitato nel tempo: la lista dei peer è
letta una sola volta all'avvio (Sezione~\ref{sec:details}C) e sia
\texttt{fetch\_peers} sia \texttt{deregister\_worker} degradano senza
propagare l'errore se il discovery server non risponde. Un guasto del
discovery server dopo l'avvio blocca quindi solo l'ingresso di nuovi client,
non il gossip già in corso tra quelli attivi, che continuano a operare sulla
propria cache locale.

\begin{figure}[t]
  \centering
  \begin{tikzpicture}[
      node distance=0.9cm and 1.1cm,
      box/.style={draw, rounded corners, minimum width=2.1cm, minimum height=0.8cm, align=center, font=\footnotesize},
      >=Stealth
    ]
    \node[box] (reg) {Discovery Server\\(REST, solo coordinamento)};
    \node[box, below left=of reg] (w0) {Worker 0};
    \node[box, below=of reg] (w1) {Worker 1};
    \node[box, below right=of reg] (w2) {Worker 2};

    \draw[dashed, ->] (w0) -- node[left, font=\tiny]{register/peers} (reg);
    \draw[dashed, ->] (w1) -- (reg);
    \draw[dashed, ->] (w2) -- node[right, font=\tiny]{register/peers} (reg);

    \draw[<->, thick] (w0) -- node[below, font=\tiny]{gossip push (gRPC)} (w1);
    \draw[<->, thick] (w1) -- (w2);
    \draw[<->, thick] (w0) to[bend right=25] (w2);
  \end{tikzpicture}
  \caption{Architettura logica: il discovery server gestisce solo la
  registrazione (linee tratteggiate REST); i modelli viaggiano esclusivamente
  peer-to-peer via gRPC (frecce continue). Il grafo tra worker è
  \emph{potenziale}, non fisso: a ogni round ciascun worker sceglie
  casualmente $k$ destinatari tra i peer disponibili, senza inoltro da parte
  di chi riceve (Sezione~\ref{sec:background}E).}
  \label{fig:architecture}
\end{figure}

\subsection{Aggregazione decentralizzata}
Ogni worker mantiene un buffer di aggregazione protetto da lock. Quando un
messaggio gossip arriva, il server gRPC non lo accoda: lo somma
immediatamente, pesato per il numero di campioni del mittente, a un
accumulatore già esistente (\emph{online aggregation}). Il costo in memoria e
in tempo di questa operazione è quindi $O(\text{dimensione del modello})$ e
non $O(K \cdot \text{dimensione del modello})$ come sarebbe mantenendo una
lista di tutti i modelli ricevuti. All'inizio del round successivo, il thread
di training applica FedAvg (Eq.~\ref{eq:fedavg}) tra il proprio modello e la
somma pesata accumulata, quindi svuota il buffer. In nessun momento un nodo
possiede o osserva i pesi di \emph{tutti} gli altri $K-1$ client
contemporaneamente: ciascun worker aggrega solo ciò che i suoi vicini gli
hanno inviato, realizzando un'aggregazione FedAvg genuinamente
peer-to-peer e priva di un aggregatore singolo per l'intera durata del
processo.

\subsection{Motivazione delle scelte protocollari}
La comunicazione tra worker usa gRPC su Protobuf anziché REST/JSON: la
serializzazione binaria di un \texttt{state\_dict} PyTorch produce un
messaggio 2--3$\times$ più compatto rispetto alla codifica testuale
equivalente, un risparmio non trascurabile quando ogni messaggio trasporta
l'intero modello (Sezione~\ref{sec:results}). La lista dei peer è letta dal
discovery server una sola volta, all'avvio (Sezione~\ref{sec:details}A), e
tenuta in una cache locale per l'intera durata del training: se il worker
la reinterrogasse a ogni round, il traffico verso il discovery server
crescerebbe linearmente con il numero di round anziché restare $O(1)$ per
l'intera run, vanificando in parte l'obiettivo di traffico ridotto che
motiva anche la scelta del k-push (Sezione~\ref{sec:background}E). La cache
viene aggiornata solo lazily, cioè non periodicamente ma esclusivamente
quando serve: ogni invio ha un timeout configurabile e, in caso di
fallimento, il worker interroga di nuovo il discovery server e ritenta una
sola volta con peer non ancora contattati in quel round — evitando sia
blocchi su nodi caduti sia un sovraccarico di richieste al discovery server
nei round in cui tutto funziona normalmente.

\subsection{Aderenza ai requisiti}
La Tabella~\ref{tab:requirements} riassume come ciascun requisito della
traccia (progetto ML e vincoli imposti per il corso di Sistemi Distribuiti e
Cloud Computing) sia stato affrontato.

\begin{table}[t]
  \centering
  \caption{Mappatura requisiti $\to$ soluzione adottata.}
  \label{tab:requirements}
  \footnotesize
  \begin{tabular}{p{2.3cm}p{5.3cm}}
    \toprule
    \textbf{Requisito} & \textbf{Come è soddisfatto} \\
    \midrule
    Team di 2 studenti & Progetto sviluppato da 2 autori; la modalità
      centralizzata con parameter server è opzionale solo per team da 3 e
      non è stata implementata. \\
    $K$ client, avvio autonomo & Ogni worker legge \texttt{config.yaml} e
      avvia il proprio ciclo di training senza comando esterno
      (Sez.~\ref{sec:design}A). \\
    Scelta della strategia di aggregazione & FedAvg selezionabile via
      \texttt{aggregation\_strategy} in configurazione; unica strategia
      implementata (Sez.~\ref{sec:background}D). \\
    Centralizzazione limitata al coordinamento & Discovery server: solo
      register/deregister/peers/health, mai pesi del modello
      (Sez.~\ref{sec:design}A). \\
    Aggregazione P2P, no aggregatore singolo & Aggregazione online nel
      buffer di ciascun worker, gossip $k$-push (Sez.~\ref{sec:design}B). \\
    LEAF, più utenti per client & FEMNIST partizionato per scrittore, blocchi
      di scrittori per worker (Sez.~\ref{sec:details}A). \\
    Scalabilità: accuracy e tempi vs $N$ & $N \in \{3,5,8\}$ testati sia in
      locale sia su AWS multi-istanza (Sez.~\ref{sec:results}). \\
    Efficienza di comunicazione & Dimensione messaggio, traffico stimato e
      tempo di Fase C misurati separatamente dal training
      (Sez.~\ref{sec:results}). \\
    Fault tolerance (opzionale) & Crash/drop injection, timeout, retry,
      staleness check (Sez.~\ref{sec:details}C). \\
    \bottomrule
  \end{tabular}
\end{table}
